\begin{lemma}[Projective effacements by passage to the opposite category]
\label{lemma-projective-effacement-opposite-category}
\begin{reference}
\cite[Chapter I, \S 1.10, Remark 4, p. 137]{Tohoku}
\end{reference}
Let $\mathcal A$ be an abelian category.
\begin{enumerate}
\item An object $P$ is projective in $\mathcal A$ if and only if it is
injective as an object of $\mathcal A^{opp}$.
\item A surjection $p:M\to A$ is a projective effacement of $A$ in
$\mathcal A$ if and only if the injection $p^{opp}:A\to M$ is an injective
effacement of $A$ in $\mathcal A^{opp}$.
\item The category $\mathcal A$ has enough projectives if and only if
$\mathcal A^{opp}$ has enough injectives.
\item If $\mathcal A^{opp}$ is a Grothendieck abelian category, then there
are a functor $P:\mathcal A\to\mathcal A$ and a natural transformation
\[
P\longrightarrow\operatorname{id}_{\mathcal A}
\]
whose component $P(A)\to A$ is a surjection from a projective object for
every $A\in\mathcal A$. In particular, these components are functorial
projective effacements.
\end{enumerate}
Thus statements about injective objects, injective effacements, and enough
injectives have the corresponding projective statements after passage to
the opposite category.
\end{lemma}

\begin{proof}
The opposite of an abelian category is abelian; monomorphisms and
epimorphisms, kernels and cokernels, and all commutative diagrams are
interchanged. Reversing the arrows in the lifting property for an injective
object gives exactly the lifting property for a projective object. Reversing
the arrows in the defining square for an injective effacement gives exactly
the defining square for a projective effacement. This proves (1) and (2),
and (3) follows by reversing an injective embedding of each object.

For (4), apply Injectives, Theorem
\ref{injectives-theorem-injective-embedding-grothendieck} to
$\mathcal A^{opp}$. Its functorial injective embeddings become, after
reversing arrows, the asserted functorial surjections in $\mathcal A$.
Their sources are projective by (1), and Lemma
\ref{lemma-projective-effacement} shows that every such surjection
is a projective effacement.
\end{proof}
